% Generated from locale/tr/content/many-valued-logic/three-valued-logics/multiple-designation.tex; do not hand-edit.


\olfileid[tr]{mvl}{thr}{mul}

\olsection{Yalnızca $\True$'nun Değil, Başka Değerlerin de Belirlenmesi}

Şimdiye kadar gördüğümüz mantıkların hepsinde belirlenmiş doğruluk değerleri
kümesi $V^+ = \{\True\}$ idi; yani bir şey ancak ve ancak doğruluk değeri
$\True$ ise doğru sayılıyordu. Fakat bir şey yalnızca $\False$ değilse de
doğru sayılabilir. O zaman örneğin matriste
$V^+ = \{\True, \Undef\}$ koşulu konarak bir mantık elde edilir.

\begin{defn}
\emph{Paradoks mantığı}~$\LogLP$ şu matrisle tanımlanır:
\begin{enumerate}
  \item $\lnot$, $\land$, $\lor$, $\lif$ bağlaçlarını içeren standart
  önerme dili $\Lang L_0$.
  \item Doğruluk değerleri kümesi $V = \{\True, \Undef, \False\}$.
  \item $\True$ ile $\Undef$ belirlenmiştir; yani
  $V^+ = \{\True, \Undef\}$.
  \item Doğruluk fonksiyonları güçlü Kleene mantığındakilerle aynıdır.
\end{enumerate}
\end{defn}

\begin{defn}
Halld\'en'in \emph{anlamsızlık mantığı}~$\LogHal$ şu matrisle tanımlanır:
\begin{enumerate}
  \item $\lnot$, $\land$, $\lor$, $\lif$ ve bir-konumlu $+$ bağlacını
  içeren standart önerme dili $\Lang L_0$.
  \item Doğruluk değerleri kümesi $V = \{\True, \Undef, \False\}$.
  \item $\True$ ile $\Undef$ belirlenmiştir; yani
  $V^+ = \{\True, \Undef\}$.
  \item Doğruluk fonksiyonları zayıf Kleene mantığındakilerle aynıdır;
  bunlara “anlamsızdır” işlemcisi eklenir:
  \begin{center}
    \begin{tabular}{c|c} 
    $\tf{+}$ & \\ 
    \hline  
    $\True$ & $\False$ \\ 
    $\Undef$ & $\True$ \\
    $\False$ & $\False$ 
    \end{tabular}
  \end{center}
\end{enumerate}
\end{defn}

Doğruluk tablolarını paylaştıkları Kleene mantıklarının tersine bunların
\emph{totolojileri vardır}.

\begin{prop}\ollabel{prop:LP-taut-CL} $\LogLP$'nin totolojileri klasik
  önerme mantığının totolojileriyle aynıdır.
\end{prop}

\begin{proof}
  \olref[syn][sub]{prop:mvl-cl} uyarınca $\Entails[\LogLP] !A$ ise
  $\Entails[\LogCL] !A$. Ters yönü göstermek için
  $\pValue v(!A)[\LogKs] = \False$ olacak biçimde
  !!a{valuation}~$\pAssign v\colon \PVar \to \{\False, \True,
  \Undef\}$ varsa $\pValue {v'}(!A)[\LogCL] = \False$ olacak biçimde
  !!a{valuation}~$\pAssign {v'}\colon \PVar \to \{\False, \True\}$
  bulunduğunu göstereceğiz. Bu, $\LogLP$ için sonucu verir; çünkü $\LogKs$
  ile $\LogLP$ aynı karakteristik doğruluk fonksiyonlarına sahiptir ve
  $\False$, $\LogLP$'nin belirlenmemiş olan tek doğruluk değeridir
  ($\LogLP$ ile~$\LogKs$ arasındaki tek fark budur). Dolayısıyla
  $\Entails/[\LogLP] !A$ ise bazı !!{valuation}~$\pAssign v$ için
  $\pValue v(!A)[\LogLP] = \pValue v(!A)[\LogKs] = \False$ olur.
  İspatlamakta olduğumuz sav uyarınca
  $\pValue{v'}[\LogCL](!A) = \False$; yani
  $\Entails/[\LogCL] !A$.
  
  Savı ispatlamak için önce $\pAssign {v'}$'yi şöyle tanımlayalım:
  \[
    \pAssign {v'}(p) = 
    \begin{cases}
    \True & \text{$\pAssign {v}(p) \in \{\True, \Undef\}$ ise}\\
    \False & \text{aksi durumda}
  \end{cases}
  \]
  Şimdi $!A$ üzerinde tümevarımla şunları göstereceğiz: (a)~Eğer
  $\pValue v(!A)[\LogKs] = \False$ ise
  $\pValue {v'}(!A)[\LogCL] = \False$; (b)~eğer
  $\pValue v(!A)[\LogKs] = \True$ ise
  $\pValue {v'}(!A)[\LogCL] = \True$.
  \begin{enumerate}
    \item Tümevarım temeli: $!A \ident p$.
    \olref[syn][val]{defn:pValue} uyarınca
    $\pValue v(!A)[\LogKs] = \pAssign v(p)$ ve
    $\pValue {v'}(!A)[\LogCL] = \pAssign {v'}(p)$. İlk değer $\False$ ise
    $\pAssign {v'}$ tanımı gereği ikinci değer de $\False$; ilk değer
    $\True$ ise ikinci değer de $\True$ olur. Böylece hem (a) hem de (b)
    elde edilir.
    
    Tümevarım adımı için şu durumları ele alın:
    \item $!A \ident \lnot !B$. 
    \begin{enumerate}
      \item $\pValue v(\lnot!B)[\LogKs] = \False$ olduğunu varsayın.
      $\tf{\lnot}[\LogKs]$ tanımı gereği
      $\pValue v(!B)[\LogKs] = \True$. Tümevarım varsayımının (b) durumundan
      $\pValue {v'}(!B)[\LogCL] = \True$ elde edilir; dolayısıyla
      $\pValue {v'}(\lnot !B)[\LogCL] = \False$.
      \item $\pValue v(\lnot!B)[\LogKs] = \True$ olduğunu varsayın.
      $\tf{\lnot}[\LogKs]$ tanımı gereği
      $\pValue v(!B)[\LogKs] = \False$. Tümevarım varsayımının (a) durumundan
      $\pValue {v'}(!B)[\LogCL] = \False$ elde edilir; dolayısıyla
      $\pValue {v'}(\lnot !B)[\LogCL] = \True$.
    \end{enumerate}

    \item $!A \ident (!B \land !C)$.
    \begin{enumerate}
      \item $\pValue v(!B \land !C)[\LogKs] = \False$ olduğunu varsayın.
      $\tf{\land}[\LogKs]$ tanımı gereği
      $\pValue v(!B)[\LogKs] = \False$ ya da
      $\pValue v(!C)[\LogKs] = \False$. Tümevarım varsayımının (a)
      durumundan $\pValue {v'}(!B)[\LogCL] = \False$ ya da
      $\pValue {v'}(!C)[\LogCL] = \False$ elde edilir; dolayısıyla
      $\pValue {v'}(!B \land !C)[\LogCL] = \False$.
      \item $\pValue v(!B \land !C)[\LogKs] = \True$ olduğunu varsayın.
      $\tf{\land}[\LogKs]$ tanımı gereği
      $\pValue v(!B)[\LogKs] = \True$ ve
      $\pValue v(!C)[\LogKs] = \True$. Tümevarım varsayımının (b)
      durumundan $\pValue {v'}(!B)[\LogCL] = \True$ ve
      $\pValue {v'}(!C)[\LogCL] = \True$ elde edilir; dolayısıyla
      $\pValue {v'}(!B \land !C)[\LogCL] = \True$.
    \end{enumerate}
  \end{enumerate}

  Öteki iki durum benzerdir ve alıştırma olarak bırakılmıştır. Başka bir
  yoldan, yukarıdaki ispat sonucu yalnızca $\lnot$ ile~$\land$ içeren bütün
  !!{formula}s için kurar. Şimdi hem $\LogKs$ hem de $\LogCL$ içinde her
  $\pAssign v$ için $\pValue v(!B \lor !C) =
  \pValue v(\lnot(\lnot!B \land \lnot !C))$ ve
  $\pValue v(!B \lif !C) = \pValue v(\lnot(!B \land \lnot !C))$
  olgularına başvurulabilir.
\end{proof}

\begin{prob}
  \olref[mvl][thr][mul]{prop:LP-taut-CL} ispatını tamamlayın; yani
  $!A \ident (!B \lor !C)$ ve $!A \ident (!B \lif !C)$ durumları için
  (a) ile~(b)'yi kurun.
\end{prob}

\begin{prob}
Her klasik totolojinin $\LogHal$ içinde bir totoloji olduğunu ispatlayın.
\end{prob}

Klasik mantıkla aynı totolojilere sahip olsalar da gerektirme bağıntıları
farklıdır. Örneğin $\LogLP$, $\lnot p, p \Entails/ q$ olması anlamında
\emph{paratutarlıdır}; dolayısıyla patlama ilkesi
$\lnot !A, !A \Entails !B$ genel olarak geçerli değildir. (Bazı $!A$ ve
$!B$ durumlarında, örneğin $!B$ bir totolojiyse geçerlidir.)

\begin{prob}
  Aşağıdaki bağıntılardan hangileri (a)~$\LogLP$ ve
  (b)~$\LogHal$ içinde geçerlidir? Her biri için bir doğruluk tablosu verin.
  \begin{enumerate}
    \item $p, p \lif q \Entails q$
    \item $\lnot q, p \lif q \Entails \lnot p$
    \item $p \lor q, \lnot p \Entails q$
    \item $\lnot p, p \Entails q$
    \item $p \Entails p \lor q$
    \item $p \lif q, q\lif r \Entails p \lif r$
  \end{enumerate}
\end{prob}

$\LogLuk[3]$ içinde $\Undef$'ı belirlenmiş sayarsanız ne olur?

\begin{defn}
  \emph{3-değerli R-Mingle} mantığı~$\LogRM[3]$ şu matrisle tanımlanır:
  \begin{enumerate}
    \item $\lfalse$, $\lnot$, $\land$, $\lor$, $\lif$ bağlaçlarını içeren
    standart önerme dili $\Lang L_0$.
    \item Doğruluk değerleri kümesi $V = \{\True, \Undef, \False\}$.
    \item $\True$ ile $\Undef$ belirlenmiştir; yani
    $V^+ = \{\True, \Undef\}$.
    \item Doğruluk fonksiyonları \L ukasiewicz mantığı~$\LogLuk[3]$ ile
    aynıdır.
  \end{enumerate}
\end{defn}

\begin{prob}
  Aşağıdaki bağıntılardan hangileri $\LogRM[3]$ içinde geçerlidir?
  \begin{enumerate}
    \item $p, p \lif q \Entails q$
    \item $p \lor q, \lnot p \Entails q$
    \item $\lnot p, p \Entails q$
    \item $p \Entails p \lor q$
  \end{enumerate}
\end{prob}

Farklı doğruluk tabloları bazen yalnızca belirlenmiş değerler değiştirilerek
aynı mantığı (gerektirme bağıntısını) üretebilir. Örneğin G\"odel mantığında
$V^+ = \{\True\}$ yerine $V^+ = \{\True, \Undef\}$ alındığında böyle olur.

\begin{prop}\ollabel{prop:gl-udes}
  $V = \{\False, \Undef,\True\}$, $V^+=\{\True, \Undef\}$ ve $3$-değerli
  G\"odel mantığının doğruluk fonksiyonlarından oluşan matris klasik mantığı
  tanımlar.
\end{prop}

\begin{proof}
  Alıştırma.
\end{proof}

\begin{prob}
  \olref[mvl][thr][mul]{prop:gl-udes} önermesini şöyle ispatlayın:
  G\"odel mantığı gibi, fakat $V^+=\{\True,\Undef\}$ ile tanımlanan
  $\Log L$ mantığı için $\Gamma \Entails/[\Log L] !B$ ise
  $\Gamma \Entails/[\LogCL] !B$ olduğunu gösterin.
  \olref[mvl][thr][mul]{prop:LP-taut-CL} düşüncelerini kullanın; fakat
  (a) ve~(b) özelliklerini ispatlamak yerine
  $\pValue v(!A)[\LogGod] = \False$ ancak ve ancak
  $\pValue {v'}(!A)[\LogCL] = \False$ olduğunu (dolayısıyla
  $\pValue v(!A)[\LogGod] \in \{\True,\Undef\}$ ancak ve ancak
  $\pValue {v'}(!A)[\LogCL] = \True$ olduğunu) gösterin. Bunun önermeyi
  neden kurduğunu açıklayın.
\end{prob}

